\documentclass[10pt,a4paper]{amsart}
\usepackage[margin=19mm]{geometry}
\usepackage{amsmath,amssymb,mathtools}
\usepackage[scr=rsfs,cal=euler]{mathalfa}
\usepackage{xfrac}
\usepackage[unicode,hidelinks,bookmarks=false]{hyperref}
\newcommand{\ie}{i.e.\ }
\newcommand{\cf}{cf.\ }
\newcommand{\resp}{respectively\ }
\newcommand{\Subsection}{\subsection}
\providecommand{\nobreakdash}{\nobreak\hbox{-}}
\setlength{\emergencystretch}{2em}
\multlinegap0pt
\newcommand{\Z}{\mathbb Z}
\usepackage{amssymb}
\usepackage{mathtools}
\usepackage[shortlabels]{enumitem}
\usepackage{float}
\newtheorem{proposition}{Proposition}
\title{On a problem of Caro on the $\Z_3$-Ramsey number of forests}
\author{Jos\'e D. Alvarado, Lucas Colucci, Roberto Parente}
\date{Source: arXiv:2503.01032v3 (CC BY 4.0)}
\begin{document}
\maketitle

\noindent\emph{Source and licence.} On a problem of Caro on the $\Z_3$-Ramsey number of forests. Version arXiv:2503.01032v3 (CC BY 4.0), distributed by arXiv under the Creative Commons Attribution 4.0 International licence, http://creativecommons.org/licenses/by/4.0/, as recorded in the arXiv licence metadata for this exact version and shown on the abstract page https://arxiv.org/abs/2503.01032v3. Full licence notice: \texttt{licenses/arxiv-2503.01032-CC-BY-4.0.txt}.\par\smallskip

\noindent\textit{Editorial scope.} Counted unit: the article's proposition prop:siblings-upper (source0.tex lines 291--293). The article's complete original proof of it is delivered with it (source0.tex lines 295--346, one \texttt{proof} environment of 52 lines). No mathematics is written, repaired or re-derived: the statement and the proof are the article's own lines, in the article's own notation, and no retained line is substituted or re-flowed.  Retained uncounted as the source-local context the proof needs: lines 370--372.  Nothing was omitted from any retained span, so the package carries no omission record.  No retained line contains a citation, so the wrapper carries no bibliography and no bibliography entry.  No retained line contains a figure environment, an image-inclusion command or drawing code, so no image is delivered and the package declares no asset dependency.  Editorial formatting only, in addition: an AMS \texttt{amsart} preamble that restates the article's own declarations and macros used by the retained lines, namely the environments \texttt{proposition} on separate counters, together with the standard packages those lines need.  The retained environments print in this document as Proposition 1 in order of appearance, whereas the article prints the same items with the numbering of its own full document.  The complete native source of the article as distributed in the arXiv e-print for this exact version is bundled unchanged as \texttt{sources/174115/source0.tex}.
\begin{proposition}[zero-sum copies in CC colorings]\label{prop:cc-copy}
Let $F$ be a type $0$ forest on $m$ vertices with $3\mid e(F)$. Every CC coloring of $K_m$ contains a zero-sum copy of $F$.
\end{proposition}

\begin{proposition}\label{prop:siblings-upper}
    Let $F$ be a forest on $n \geq 7$ vertices which is not a star and not $1$ modulo $3$ regular such that $3 \mid e(F)$ and that contains two siblings leaves. Then, $R(F,\mathbb{Z}_3) \leq n+1$. 
\end{proposition}

\begin{proof}
    Let $f_1$ and $f_2$ be two leaves in $F$ with a common parent $p$.
    First of all, notice the result is immediate if the coloring $\chi$ of the complete graph $K$ contains a trichromatic vertex: indeed, if $v$ has three neighbors $u_0$, $u_1$, $u_2$ such that the color of $vu_i$ is $i$ for $i \in \{0,1,2\}$, then we may embed $F-\{f_1,f_2\}$ in $K-\{u_0,u_1,u_2\}$ sending $p$ to $v$, but otherwise arbitrarily. Then, we can embed the two remaining leaves in three different ways, choosing with two of the $u_i$ will be the image of the $f_i$. In this manner, we create three copies of $F$ in $K$ with distinct sums, so that one of them has zero sum.

    In what follows, we may assume that the coloring of $K$ does not contain a trichromatic vertex, i.e., every vertex is incident to edges of at most two different colors. This partition the vertices of $K$ into three sets: $V_{01}$, $V_{02}$ and $V_{12}$, where the vertices in the class $V_{ij}$ only have edges of colors $i$ and $j$ incident to them. The rest of the proof will be divided into cases according to the sizes of the sets $V_{ij}$:

    \textbf{Case 1:} All the $V_{ij}$ are nonempty

    \medskip
    
    Without loss of generality, we may assume that $|V_{12}| \geq 3$. Let $l_1$ and $l_2$ be leaves of $F$ with distinct parents $p_1$ and $p_2$, respectively. We embed $l_1$, $p_2$ and $l_2$ in $V_{12}$.

    If the edge $p_2l_2$ is colored with $1$ (resp. $2$), we embed $p_1$ in $V_{01}$ (resp. $V_{02}$) and leave one of the vertices of $V_{02}$ (resp. $V_{01}$) spare, embedding the rest of the forest arbitrarily. If the sum of $F$ is not zero, we can turn it to zero by moving the embedding of one of the $l_i$ to the spare vertex.

    \textbf{Case 2:} Exactly one the $V_{ij}$ is empty

    Assume that the nonempty sets are $V_{01}$ and $V_{02}$. First, notice that we may suppose that there is an edge of color $1$ in the set $V_{01}$ and an edge of color $2$ in the set $V_{02}$, for otherwise we could merge the two sets in a single $V_{ij}$, which leads us to Case 3. Furthermore, we may assume that there is an edge of color $0$ inside one of the sets, for otherwise we would have a CC coloring and Proposition \ref{prop:cc-copy} would apply. Without loss of generality, assume that it is inside $V_{02}$. In particular, there are three vertices $x$, $y$ and $z$ in $V_{02}$ such that $\chi(xy)=2$ and $\chi(yz)=0$. Let $l_1$ and $l_2$ be leaves of $F$ with distinct parents $p_1$ and $p_2$, respectively. We embed $l_1$ in $x$, $p_1$ in $y$, leave $z$ spare, and embed $l_2$ and $p_2$ in $V_{01}$ in a way that the edge $l_2p_2$ gets color $1$. Again, if the sum of $F$ is not zero, we can turn it to zero by adding or subtracting $1$ by moving one of the $l_i$ to the spare vertex.

    \medskip

    \textbf{Case 3:} Only one of the $V_{ij}$ is nonempty

    Assume without loss of generality that $V_{01}$ is the only nonempty set, i.e., all the edges of the complete graph are colored with either color $0$ or $1$.

    Suppose first that there is an alternating path on $5$ vertices in the coloring, i.e., five vertices $v_1, \dots, v_5$ such that $\chi(v_1v_2)=1$, $\chi(v_2v_3)=0$, $\chi(v_3v_4)=1$ and $\chi(v_4v_5)=0$. Let $l_1$ and $l_2$ be leaves of $F$ with distinct parents $p_1$ and $p_2$, respectively. We embed $l_1$ in $v_1$, $p_1$ in $v_2$, leave $v_3$ spare, embed $p_2$ in $v_4$, and $l_2$ in $v_5$. If the sum of $F$ is not zero, we can turn it to zero by adding or subtracting $1$ by moving one of the $l_i$ to the spare vertex.

    On other hand, if the coloring does not contain an alternating path on $5$ vertices, let us consider a longest alternating path in it. If it has just one edge, the coloring in monochromatic and the result follows immediately.

    If a longest alternating path has two edges, let $uvw$ be one alternating path with $\chi(uv)=0$ and $\chi(vw)=1$. It follows immediately that $\chi(xu)=0$ and $\chi(xw)=1$ for every vertex $x \in K-\{u,v,w\}$. Let us assume without loss of generality that $\chi(uw)=0$. Then, if $x$ and $y$ are vertices of $K-\{u,v,w\}$, considering the path $yxwu$, it follows that $\chi(xy)=1$, i.e., $K-\{u,v,w\}$ is monochromatic of color $1$. Finally, considering the path $xvwu$, it follows that $\chi(vx)=1$ for every $x \in K-\{u,v,w\}$. This implies that $K-\{u\}$ is monochromatic, so there is a zero-sum copy of $F$ in it. 

    Finally, suppose that a longest alternating path has three edges. Let $C$ be a maximum monochromatic clique in $K$. It has at least three vertices (as $R(3,3)=6$). Let us assume without loss of generality that it has color $0$. 

    First, note that $K-C$ is monochromatic. Indeed, if it is not the case, it contains three vertices $x$, $y$ and $z$ with $\chi(xy)=1$ and $\chi(yz)=0$. Also, as $C$ is maximal, there is $u \in C$ such that $\chi(zu)=1$. Let $t$ be another vertex in $C$. Then, the path $xyzut$ is alternating, a contradiction.

    If $K-C$ is monochromatic of color $0$, then we cannot have two independent edges of color $1$ in $K$: indeed, if $\chi(xy)=\chi(zw)=1$ where $x, w \in C$, $y, z \in K-C$ are four distinct vertices, let $t$ be another vertex of $C$. Then, $xyzwt$ is an alternating path on $4$ edges, a contradiction. This proves that all the edges of color $1$ are incident to a single vertex. Deleting this vertex, we get a monochromatic complete graph on $n$ vertices, so there is a zero-sum copy of $F$ in it.

    Finally, if $K-C$ is monochromatic of color $1$, we claim that one of the two colors spans a clique.
    
    
    For that purpose, note that for every $y \in C$ all the edges joining $y$ and $K-C$ have the same color. Indeed, if it is not the case, let $x, z \in K-C$, $y \in C$ such that $\chi(xy)=1$ and $\chi(yz)=0$. By the maximality of $C$, there is a vertex $w \in C$ with $\chi(zw)=1$. Let $t$ be another vertex of $C$. Then, the path $xyzwt$ is alternating, a contradiction.

    If there are two vertices $y_1$ and $y_2 \in C$ such that $c(xy_1)=c(xy_2)=1$ for every $x \in K-C$ and a vertex $y_3 \in C$ such that $c(xy_3)=0$ for every $x \in K-C$, then the path $xy_1y_2wy_3$ is alternating for every $x, w \in K-C$, a contradiction.

    This implies that either all the edges in between $C$ and $K_C$ are colored $1$, or there is a vertex $y \in C$ such that $K-C\cup \{y \}$ is the clique spanned by color $1$, which proves the claim.

    Suppose without loss of generality that the clique $Q$ spanned by one of the colors is monochromatic with color $0$ and all other edges in the graph have color $1$.

    If  $|K-Q| \leq 1$ or $|Q| \leq 2$, then $K$ contains a monochromatic $K_n$, and hence a zero-sum copy of $F$. Also, if $|K-Q|=2$, say, $K-Q = \{u,v\}$ then we get a zero-sum copy of $F$ by embedding either a vertex of degree zero modulo $3$ of $F$ in $u$ and leaving $v$ spare, or embedding a vertex $x$ of degree $2$ modulo $3$ in $u$ and a leaf that is not adjacent to $x$ in $v$.

    From now on, we may assume that $|K-Q| \geq 3$ and $|Q| \geq 3$. Let $p$ be a parent of a leaf $f$ with $2 \leq d(p) \leq n-3$ (such a vertex exists from the hypothesis of the statement). We will embed $F$ in $K$ in a way such that $p$ and a number congruent to $2$ modulo $3$ of its neighbors, including $f$, are embedded in $Q$, and the spare vertex of $K$ is in $K-Q$. If we have such an embedding, by switching either $f$ or $p$ with the spare vertex, we can get copies of $F$ with zero sum. To get such an embedding, we first embed $p$, $f$ and as many neighbors of $p$ as possible in $Q$. Moving at most two neighbors of $p$ to $K-Q$, we get $2$ modulo $3$ neighbors in $Q$ and a spare vertex in $K-Q$.
    
\end{proof}

\end{document}
